\section{Portfolio Construction Model}

The optimizer is chosen only after the cashflow and risk map are built. Covariance is not
enough. A high Sharpe ratio cannot come from \(\Sigma_{O,t}\) alone; it must come from a
conditional expected-return vector. The paper writes the option mean as
\[
\widehat\mu_{i,t}
=
 b^S_{i,t}\widehat\lambda^S_t
+b^\Gamma_{i,t}\widehat\lambda^{var}_t
+b^\nu_{i,t}\widehat\lambda^{vol}_t
+b^{skew}_{i,t}\widehat\lambda^{skew}_t
+b^{tail}_{i,t}\widehat\lambda^{tail}_t
+\theta^\ast_{i,t}.
\]
Delta loads on equity compensation. Gamma and vega load on variance and volatility risk
premia. Put skew and VIX call wings load on crash-insurance and tail premia. Theta is
carry under unchanged state variables. The option-return literature documents why these
components can carry expected returns rather than only hedge noise
\citep{coval_shumway2001,bakshi_kapadia2003,carr_wu2009,
bollerslev_tauchen_zhou2009,broadie_chernov_johannes2009}. In implementation the
combined signal is estimated only from training-window or decision-date information and is
shrunk toward zero.

Given \(\widehat\mu_t\) and \(\widehat\Sigma_t\), the gross option-only Sharpe problem is
\[
  \max_{q\in\mathcal K_t}
  \frac{q^\top\widehat\mu_t}{\sqrt{q^\top\widehat\Sigma_tq}},
  \qquad
  \widehat\Sigma_t=B_t\Omega_tB_t^\top+\Sigma_{\varepsilon,t}.
\]
If the feasible set is unconstrained and \(\widehat\Sigma_t\) is positive definite, the
maximum-Sharpe direction is proportional to
\(\widehat\Sigma_t^{-1}\widehat\mu_t\). This is exactly the Markowitz tangency formula. The
option work is in building the conditional option return object that the formula consumes.

The practical feasible set is smaller. The research implementation imposes gross NAV, short
premium, per-contract, underlying concentration, delta, SPY beta, eight-underlying beta,
gamma, equity vega, VIX vega, and scenario stress-loss budgets where the corresponding
strategy requires them. The optimizer can therefore choose calls, puts, long positions, and
short positions, but it cannot hide large exposure behind a small option premium.

\begin{table}[H]
\centering
\scriptsize
\IfFileExists{tables/risk_calibration.tex}
  {\resizebox{0.94\textwidth}{!}{\begin{tabular}{lrrrr}
\toprule
Strategy & Predicted monthly vol & Train realized vol & Test realized vol & Test / predicted \\
\midrule
Equity-option Greek Markowitz & 0.334 & 0.267 & 0.307 & 0.918 \\
Greek Markowitz + VIX & 0.363 & 0.261 & 0.262 & 0.722 \\
Beta/delta-neutral + VIX & 0.354 & 0.272 & 0.239 & 0.677 \\
Equal premium & 1.509 & 0.625 & 0.542 & 0.359 \\
Equal risk & 0.668 & 0.394 & 0.368 & 0.551 \\
VIX hedge sleeve & 4.648 & 4.092 & 0.417 & 0.090 \\
\bottomrule
\end{tabular}
}}
  {\begin{tabular}{llll}\toprule Strategy & Predicted risk & Realized risk & Ratio \\\midrule
  Missing & Run the empirical pipeline & -- & -- \\\bottomrule\end{tabular}}
\caption{Greek-induced covariance calibration. Predicted monthly volatility is computed
from \(q^\top\Sigma_Oq\) using the training-window Greek risk model; realized volatility is
measured on train and test option returns.}
\label{tab:risk_calibration}
\end{table}

The estimator follows the standard risk-model route. It constructs \(B_t\), estimates
\(\Omega_t\) from point-in-time state shocks, estimates \(\Sigma_{\varepsilon,t}\) from
training residuals, shrinks the covariance toward a diagonal target, and repairs small
negative eigenvalues. This follows the same concern that motivates covariance shrinkage in
ordinary portfolio problems: noisy matrices become dangerous when the optimizer inverts
them \citep{ledoitwolf2004}.
